\documentclass{situationist-preprint}

\title{Monism Without the Argument From Entanglement}
\msid{MET-09-b}
\msfamily{MET-09}
\msdiscipline{Metaphysics}
\mstype{Alternative proof}
\mskeywords{priority monism; grounding; totality facts; environment; entanglement; common fate}

\begin{document}
\maketitle

\begin{abstract}
MET-09-a argues that the cosmos is the only fundamental object. Its main argument runs through the laws: composition depends on dynamics, dynamics on which generalizations are laws, and lawhood on a totality fact about the cosmos; a second argument, from quantum entanglement, covers the natures of simples. We give a route to the same conclusion that uses neither the laws nor entanglement. Composition depends on the regime of influences that a plurality's environment exerts, and every full ground of a fact about that regime includes a totality fact. The argument therefore holds for a monist who rejects the best-system account of laws and who declines to rest metaphysics on the interpretation of quantum mechanics.
\end{abstract}

\section{What is avoided}

The main argument of MET-09-a (Section~4) uses five premises in addition to the grounding framework: Existence first, Composition grounds, Dynamics and Cosmic totality (MET-09-a, P1--P4), and No simples (Assumption~3.7). Dynamics links composition to lawhood, and Lemma~4.1 of MET-09-a then derives the need for a totality fact from the globality of lawhood established by MET-12. The argument from entanglement (MET-09-a, Section~5) adds Non-separability (Assumption~5.2) and Determination (P5).

The proof below avoids Dynamics, Lemma~4.1 and with them every appeal to the theory of laws, and it avoids Section~5 entirely. It keeps Existence first, Composition grounds, Cosmic totality, No simples, and the assumptions that partial grounding is a strict partial order (MET-09-a, Assumption~3.5) and that full grounds necessitate (Assumption~3.6). No simples is supported by the case for divisibility in MET-10-a (Section~4.1), which rests on the continuity of physical fields and not on entanglement.

\section{The environmental route}

In the criterion of MET-06-a, the members of a plurality share a common fate if, under every admissible perturbation, every two of them have the same fate; and a perturbation is admissible if its magnitude falls within the regime of external influences to which the plurality's environment characteristically exposes it (MET-06-a, Definitions~3.2--3.4). Call the fact that the environment of the $xx$ exerts on them a regime $R$ of influences the \emph{regime fact} for the $xx$. Whether the $xx$ share a common fate depends on which perturbations are admissible, and so on the regime fact.

\begin{principle}[Regime]\label{p:regime}
For every plurality $xx$ of proper parts of the cosmos, the fact that the $xx$ share a common fate under their relevant dynamics, or fail to, is partly grounded in the regime fact for the $xx$.
\end{principle}

The Regime principle replaces Dynamics. It does not say that the regime fact is grounded in laws, or that the dynamics of the $xx$ are; it says only that what counts as an admissible perturbation of the $xx$, and hence whether they hold together under every admissible perturbation, depends on what their environment does to them. A pebble on a beach and a pebble in a rock tumbler have the same internal bonds and different regimes, and whether a loosely cemented aggregate composes a stone can differ between them.

\begin{lemma}\label{l:regime}
Every full ground of a regime fact includes a totality fact.
\end{lemma}

\begin{proof}
Let $\rho$ be the regime fact for the $xx$, saying that the influences their environment exerts on them are of regime $R$, and let $\Gamma$ be a full ground of $\rho$. Suppose $\Gamma$ contains no totality fact, so that it consists of facts about which objects exist in the environment and what they are like, none of which says that these are all the environing objects. Then there is a possible world in which the members of $\Gamma$ obtain and the environment contains, in addition, a further object exerting influences outside $R$: a massive body nearby, or a source of heat. In that world the regime is not $R$, so $\rho$ does not obtain, contrary to Assumption~3.6 of MET-09-a. Hence $\Gamma$ includes a totality fact.
\end{proof}

Lemma~\ref{l:regime} is the environmental analogue of Lemma~4.1 of MET-09-a, and it rests on the same structural point \citep{Armstrong1997}: positive facts about some things never settle that there is nothing more, and a regime is fixed as much by the absence of further influences as by the presence of those that are exerted. Unlike Lemma~4.1, it needs no premise about what makes a generalization a law.

\begin{lemma}\label{l:topdown}
For every proper part $p$ of the cosmos, the totality fact $T$, that the cosmos is all there is, partly grounds $E(p)$.
\end{lemma}

\begin{proof}
By No simples, $p$ is a composite; let the $xx$ be its parts at some level. By Composition grounds, $E(p)$ is fully grounded in the fact $\kappa$ that the $xx$ share a common fate, so $\kappa \prec E(p)$. By the Regime principle, the regime fact $\rho$ for the $xx$ satisfies $\rho \prec \kappa$. By Lemma~\ref{l:regime}, every full ground of $\rho$ includes a totality fact, which by Cosmic totality and the argument of MET-09-a (Section~4.2) is $T$ or is partly grounded in $T$. So $T \prec \rho$, and by transitivity $T \prec E(p)$.
\end{proof}

\begin{theorem}\label{t:main}
The cosmos is fundamental, and it is the only fundamental concrete object.
\end{theorem}

\begin{proof}
Lemma~\ref{l:topdown} is Lemma~4.2 of MET-09-a, obtained by a different route. The rest of the proof of Theorem~4.4 of MET-09-a uses only that lemma, Existence first and Assumption~3.5: if a fact about a proper part $p$ partly grounded $E(c)$, then $E(p) \prec E(c) \prec T \prec E(p)$, which irreflexivity forbids; so the cosmos depends on no proper part, and every proper part depends on the cosmos through $T$.
\end{proof}

\section{Objections}

\begin{objection}[Local regimes]
The regime of a table is fixed by the room it stands in. Nothing in Andromeda bears on it, and the totality fact is idle.
\end{objection}

\begin{reply}
That nothing in Andromeda bears on the table is part of what fixes the regime, and it is a negative fact about the whole environment, not a fact about the room. A room in a world with a further body that perturbs it is a room with a different regime. The totality fact is not idle in the ground; it is the clause that makes the local description complete. What is true is that the totality fact is seldom informative, since it is constant across the regime facts of all proper parts. That is the expected profile of a fact about the one thing every proper part has as its outside.
\end{reply}

\begin{objection}[Descent]
MET-10 holds that every level is grounded in a finer level. If the existence of every proper part is grounded in facts about the whole, the descent is cut off.
\end{objection}

\begin{reply}
The descent concerns proper parts: facts at each level are fully grounded in facts at the next finer level, without end. Lemma~\ref{l:topdown} adds that each such fact is also partly grounded in a fact about the whole, which is not a level and not a member of one. The structure is the one shown consistent in MET-10-a (Proposition~6.1) and MET-09-a (Proposition~6.1): the descent among the parts, and the priority of the whole.
\end{reply}

\section{Conclusion}

Priority monism does not depend on quantum entanglement or on the best-system account of laws. Composition is relative to the regime of influences an environment exerts, every regime fact includes the fact that there is nothing more in the environment, and that fact is a fact about the cosmos. The cosmos, the one object without an environment, is the one object that depends on nothing \citep{Schaffer2010}.

\begin{thebibliography}{9}\small
\bibitem[Armstrong(1997)]{Armstrong1997} Armstrong, D.~M. (1997). \emph{A World of States of Affairs}. Cambridge University Press.
\bibitem[Schaffer(2010)]{Schaffer2010} Schaffer, J. (2010). Monism: The priority of the whole. \emph{Philosophical Review}, 119(1), 31--76.
\end{thebibliography}

\end{document}
